\documentclass[10pt,a4paper]{amsart}
\usepackage[margin=19mm]{geometry}
\usepackage{amsmath,amssymb,mathtools}
\usepackage[scr=rsfs,cal=euler]{mathalfa}
\usepackage{xfrac}
\usepackage[unicode,hidelinks,bookmarks=false]{hyperref}
\newcommand{\ie}{i.e.\ }
\newcommand{\cf}{cf.\ }
\newcommand{\resp}{respectively\ }
\newcommand{\Subsection}{\subsection}
\providecommand{\nobreakdash}{\nobreak\hbox{-}}
\setlength{\emergencystretch}{2em}
\multlinegap0pt
\usepackage{amssymb}
\newtheorem{lemma}{Lemma}
\newcommand{\BO}{\text{\rm BO} }
\newcommand{\Gd}{\wt{\mathcal{G}}_\dl}
\newcommand{\Gdl}{\mathcal{G}_{\dl} }
\newcommand{\KDV}{\text{\rm KdV} }
\renewcommand{\L}{\mathcal{L}}
\newcommand{\R}{\mathbb{R}}
\newcommand{\dl}{\delta}
\newcommand{\dt}{\partial_t}
\newcommand{\dx}{\partial_x}
\newcommand{\ges}{\gtrsim}
\newcommand{\les}{\lesssim}
\DeclareMathOperator{\med}{med}
\newcommand{\noi}{\noindent}
\newcommand{\wt}{\widetilde}
\title{On the singular nature of shallow-water convergence of the intermediate long wave equation on the real line}
\author{Andreia Chapouto, Benjamin Harrop-Griffiths, Guopeng Li,, Tadahiro Oh}
\date{Source: arXiv:2602.20695v1 (CC BY 4.0)}
\begin{document}
\maketitle

\noindent\emph{Source and licence.} On the singular nature of shallow-water convergence of the intermediate long wave equation on the real line. Version arXiv:2602.20695v1 (CC BY 4.0), distributed by arXiv under the Creative Commons Attribution 4.0 International licence, http://creativecommons.org/licenses/by/4.0/, as recorded in the arXiv licence metadata for this exact version and shown on the abstract page https://arxiv.org/abs/2602.20695v1. Full licence notice: \texttt{licenses/arxiv-2602.20695-CC-BY-4.0.txt}.\par\smallskip

\noindent\textit{Editorial scope.} Counted unit: the article's lemma LEM:res (source0.tex lines 1470--1513). The article's complete original proof of it is delivered with it (source0.tex lines 1516--1626, one \texttt{proof} environment of 111 lines). No mathematics is written, repaired or re-derived: the statement and the proof are the article's own lines, in the article's own notation, and no retained line is substituted or re-flowed.  Retained uncounted as the source-local context the proof needs: lines 704--708; lines 1213--1216; lines 1223--1227; lines 1280--1283; lines 1296--1305; lines 1313--1318; lines 1451--1458.  Nothing was omitted from any retained span, so the package carries no omission record.  No retained line contains a citation, so the wrapper carries no bibliography and no bibliography entry.  No retained line contains a figure environment, an image-inclusion command or drawing code, so no image is delivered and the package declares no asset dependency.  Editorial formatting only, in addition: an AMS \texttt{amsart} preamble that restates the article's own declarations and macros used by the retained lines, namely the environments \texttt{lemma} on separate counters, together with the standard packages those lines need.  The retained environments print in this document as Lemma 1 in order of appearance, whereas the article prints the same items with the numbering of its own full document.  The complete native source of the article as distributed in the arXiv e-print for this exact version is bundled unchanged as \texttt{sources/195209/source0.tex}.
\begin{align}
\dt v   -    \Gd   \dx^2 v= \dx(v^2), 
\qquad \text{where }\ \Gd := \frac3\dl \Gdl.
\label{sILW1}
\end{align}

\begin{align}
\Xi_\KDV (\bar \xi) &= 3\xi \xi_1 \xi_2,\label{res1}\\
|\Xi_{\BO}(\bar \xi)| &= 2\xi_{\med}\xi_{\min}\label{res3}.
\end{align}

\begin{align}
\Xi_\dl (\bar \xi)
= \Xi_\dl(\xi,\xi_1,\xi_2) =  p_\dl(\xi)  -   p_\dl(\xi_1) -  p_\dl(\xi_2), 
\label{res4}
\end{align}

\begin{align}
\Xi_\dl(\bar \xi) =  \Xi_\BO(\bar \xi) + O(\dl^{-2}), 
\label{res8}
\end{align}

Lastly, given $0 < \dl < \infty$, we consider the resonance function 
$\wt \Xi_\dl (\bar \xi)$  for the scaled ILW~\eqref{sILW1}, given by 
\begin{align}
\begin{split}
\wt \Xi_\dl (\bar \xi)
& = \wt \Xi_\dl(\xi,\xi_1,\xi_2) 
=  \wt p_\dl(\xi)  -  \wt  p_\dl(\xi_1) -  \wt p_\dl(\xi_2), 
\end{split}
\label{res9}
\end{align}

 \begin{align}
\wt p_\dl(\xi)
&  =   \frac 3 \dl p_\dl(\xi) 
= \xi \L_\dl(\xi),  %=  \xi^3 \big(1- h(\dl, \xi)  \big), 
\label{res10}
\end{align}

\begin{align}
\wt \Xi_\dl(\bar \xi)
= 6\xi \xi_1 \xi_2
 \sum_{k=1}^{\infty}   
 \frac{ \pi^2 k^2 \big(3 \pi^2 k^2 + \dl^2(\xi_1^2 + \xi_1 \xi_2 + \xi_2^2)\big)}
 {\prod_{j = 0}^2(\pi^2 k^2 +\dl^2\xi_j^2)}
\label{LL1}
\end{align}

\begin{lemma}\label{LEM:res}
The following holds 
for any $\xi, \xi_1, \xi_2 \in \R$
with $\xi = \xi_1 + \xi_2$, 
uniformly in $0 < \dl \le 1$\textup{:}

\smallskip

\begin{itemize}
\item
[\textup{(i)}]
Suppose that  $\xi_{\max} \les \dl^{-1}$.
Then, we have 
$|\wt \Xi_\dl(\bar \xi)|
\sim |\Xi_\KDV(\bar \xi)|\sim |\xi\xi_1\xi_2|$.




\medskip

\item [\textup{(ii)}]
Suppose that $ \xi_{\max} \gg \dl^{-1}$.
Then, we have 
$|\wt \Xi_\dl(\bar \xi)|
\sim    \dl^{-1} \xi_{\min}\xi_{\max}$. 


%\medskip
%
%\noi
%\item[\textup{(iii)}]
%Suppose that 
%$\xi_{\max} \gg    \dl^{-1} \ges  \xi_{\min}$.
%Then, we have 
%$|\wt \Xi_\dl(\bar \xi)|
% \sim    \dl^{-1} \xi_{\min}\xi_{\max}$.


\end{itemize}



\end{lemma}

\begin{proof}
(i) 
Suppose that  $\xi_{\max} \les \dl^{-1}$.
Then,  from \eqref{LL1} and \eqref{res1}, we have
\begin{align*}
|\wt \Xi_\dl(\bar \xi)|
\sim  |\xi \xi_1 \xi_2|
 \sum_{k=1}^{\infty}   
 \frac1{\pi^2 k^2}
\sim  |\xi \xi_1 \xi_2|
\sim |\Xi_\KDV(\bar \xi)|, 
\end{align*}

\noi
uniformly in $0 < \dl \le 1$.






\medskip

\noi
(ii)
We first consider the case  $ \xi_{\min} \gg \dl^{-1}$.
Then,  from \eqref{res3}, 
 we have
\begin{align}
|\Xi_\BO(\bar \xi)|\sim \xi_{\min}\xi_{\max} \gg \dl^{-2}.
\label{LL3}
\end{align}

\noi
Then, 
from 
  \eqref{res8} and  \eqref{LL3}, 
we obtain
\begin{align*}
|\Xi_\dl(\bar \xi) - \Xi_\BO(\bar \xi)| = O(\dl^{-2}) \ll 
 |\Xi_\BO(\bar \xi)|, 
\end{align*}


\noi
from which we conclude that 
\begin{align}
|\Xi_\dl(\bar \xi)|\sim | \Xi_\BO(\bar \xi)|
\label{LL4}
\end{align}

\noi
uniformly in $0 < \dl \le 1$.
Thus, from 
\eqref{LL4} with 
\eqref{res9},  \eqref{res10}, 
and  \eqref{res4},
we obtain
\begin{align*}
|\wt \Xi_\dl(\bar \xi)|\sim \dl^{-1} | \Xi_\BO(\bar \xi)|
\sim    \dl^{-1} \xi_{\min}\xi_{\max}, 
\end{align*}

\noi
uniformly in $0 < \dl \le 1$.


Next, we consider the case
  $\xi_{\max} \gg    \dl^{-1} \ges  \xi_{\min}$.
Without loss of generality, 
assume that $|\xi_1|\le |\xi_2|$.
Since $\xi_{\med} \sim \xi_{\max}$ under $\xi = \xi_1 + \xi_2$, 
it follows from 
 \eqref{LL1} with 
 $\xi_1^2 + \xi_1 \xi_2 + \xi_2^2 \ge
\frac 12  \xi_1^2 + \frac 12 \xi_2^2
\sim \xi_2^2$ that 
\begin{align*}
|\wt \Xi_\dl(\bar \xi)|
\sim |\xi \xi_1 \xi_2|
 \sum_{k=1}^{\infty}   
 \frac{ \pi^2 k^2 }
 {\prod_{j = 0}^1(\pi^2k^2  +\dl^2\xi_j^2)}.
\end{align*}

\noi
Note that the expression on the right-hand side is symmetric in $\xi$ and $\xi_1$.
Without loss of generality, assume that $|\xi| = |\xi_0| \le |\xi_1|$.
Then, we have $|\xi_1| \sim \xi_{\max}
\gg \dl^{-1} \ges \xi_{\min} = |\xi|$. 
In particular, we have 
$\pi^2k^2  +\dl^2\xi^2 \sim \pi^2 k^2$.
Then, by  estimating
the sum  by an integral via a Riemann sum approximation, we have 
\begin{align*}
|\wt \Xi_\dl(\bar \xi)|
& \sim \dl^{-1} |\xi  \xi_2|
 \sum_{k=1}^{\infty}   
 \frac{ 1}
 {\pi^2(\frac k{\dl \xi_1})^2  +1} \frac{1}{\dl |\xi_1|}\\
 & \sim 
    \dl^{-1} \xi_{\min}\xi_{\max}
\int_{0}^\infty
 \frac{ 1}
 {\pi^2x^2 +1}dx\\
& \sim    \dl^{-1} \xi_{\min}\xi_{\max}, 
\end{align*}

\noi
uniformly in $0 < \dl \le 1$.
\end{proof}

\end{document}
